\documentclass[preview]{standalone}

\usepackage{amsmath}
\usepackage{amssymb}
\usepackage{stellar}
\usepackage{definitions}

\begin{document}

\id{psicologia-percezione-contesto}
\genpage

\section{Percezione, contesto e aspettative}

\begin{snippet}{contesto-aspettative-percezione}
    L'obiettivo principale della percezione è fornire una visione utile e
    affidabile del mondo. Tuttavia, quando i processi bottom-up non permettono
    un'identificazione univoca dello stimolo, i processi top-down usano
    contesto, aspettative, memoria e conoscenze pregresse per attribuire un
    significato a ciò che viene percepito.
\end{snippet}

\begin{snippetdefinition}{illusione-definition}{Illusione}
    Un'\emph{illusione} si verifica quando uno stimolo viene percepito in modo
    non corretto rispetto alle sue caratteristiche fisiche.
\end{snippetdefinition}

\begin{snippetexample}{figura-anatra-coniglio-example}{Figura anatra-coniglio}
    La figura anatra-coniglio è una figura ambigua: lo stesso stimolo visivo può
    essere organizzato percettivamente come un'anatra oppure come un coniglio.
    L'esempio mostra che la percezione non dipende solo dai dati sensoriali, ma
    anche dall'interpretazione dello stimolo.
\end{snippetexample}

\includesnpt[width=45\%,src=/snippet/static-psychology/duck-rabbit-ambiguous-figure.jpg]{centered-img}

\begin{snippetexample}{illusione-ebbinghaus-example}{Illusione di Ebbinghaus}
    Nell'\emph{illusione di Ebbinghaus}, due cerchi centrali della stessa
    dimensione vengono percepiti diversamente a seconda dei cerchi che li
    circondano. Il cerchio circondato da cerchi più piccoli appare più grande
    dell'altro. L'effetto mostra che il raggruppamento percettivo e il contesto
    influenzano il giudizio sulle dimensioni.
\end{snippetexample}

\includesnpt[width=68\%,src=/snippet/static-psychology/ebbinghaus-illusion.jpg]{centered-img}

\begin{snippet}{equilibrio-bottom-up-top-down}
    Se la percezione fosse completamente bottom-up, registreremmo solo la
    realtà contingente del momento, senza usare l'esperienza passata per
    interpretarla. Se fosse completamente top-down, rischieremmo invece di
    vedere soprattutto ciò che ci aspettiamo o desideriamo vedere. L'equilibrio
    tra i due tipi di processo permette di rispondere in modo efficace alle
    necessità biologiche e sociali dell'individuo.
\end{snippet}

\includesnpt[width=48\%,src=/snippet/static-psychology/bottom-up-top-down-perception-balance.jpg]{centered-img}

\end{document}
